\thispagestyle{empty}
\selectlanguage{hebrew}

    \textbf{תקציר}

בשנים האחרונות גברה ההתעניינות של קהילת פיזיקה של אנרגיות גבוהות והקהילה של פיזיקה גרעינית בשימוש של מאיצי  הדרונים כמאיצי פוטונים לטובת מחקר אינטראקציות פוטון-פוטון והדרון-פוטון.  המטען של החלקיק המאיץ מייצר שטף פוטונים המאפשר ביצוע של מדידות המשתמשות בפוטונים כגשוש לדינמיקה של הכוח החזק, מבנה הגרעין , ופיזיקה של רווית גלואונים.  בהצעה זו אני מציע ללמוד את התהליכי האינטרקציה של הפוטונים הללו, תוך שימוש במידע שנאסף בגלאי אטלס במאיץ החליקים ההדרוני הגדול.  במסמך זה נפרט על התאוריה של אינטראקציות פוטונים במאיץ, ואת אוסף תהליכים האולטרה פריפריאלים בהתנגשויות יונים כבדים שמאפשרים את מדידתם. לאחר מכן נפרט על שאלות פתוחות הנוגעות במבנה התנגשויות מסיום, כגון עמימות לגבי כיוון הפוטון כאשר אין שבירה של גרעינים, וכיצד להתמודד עמם בתהליכים דו סילוניים. לאחר מכן אציע שימוש בלמידת מכונה על מנת להבדיל בין התהליכים השונים תחת התנגשויות אולטרה פריפראליות, וכיצד ניתן להשתמש במערכת זו על מנת לחלץ מידע נוסף לגבי האירועים הללו.

\selectlanguage{english}